\documentclass[11pt]{article}

% Layout.
\usepackage[top=1.2in, bottom=0.9in, left=1in, right=1in, headheight=1in, headsep=6pt]{geometry}

% Fonts.
\usepackage{mathptmx}
\usepackage[scaled=1.0]{helvet}
\renewcommand{\emph}[1]{\textsf{\textbf{#1}}}

% Misc packages.
\usepackage{amsmath,amssymb,latexsym}
\usepackage{graphicx,hyperref}
\usepackage{array}
\usepackage{xcolor}
\usepackage{multicol}
\usepackage{tabularx,colortbl}
\usepackage{enumitem}
\usepackage{soul}
\usepackage{needspace}

\hypersetup{
    colorlinks=true,
    linkcolor=blue,
    filecolor=magenta,      
    urlcolor=blue,
    pdftitle={Syllabus for MATH F113X (\S 001) Fall 2026},
    }

\def\mailto#1{\href{mailto:#1}{#1}}

% Paragraph spacing
\parindent 0pt
\parskip 6pt plus 1pt
\def\tableindent{\hskip 0.5 in}
\def\ts{\hskip 1.5 em}



\usepackage{fancyhdr}
\pagestyle{fancy} 
\lhead{\large\sf\textbf{MATH F113X Math and Society (\S 001, async)}}
\rhead{\large\sf\textbf{Fall 2026 Syllabus}}

\newcommand{\localhead}[1]{\par\smallskip\textbf{#1} \smallskip\nobreak\\}%
\def\heading#1{\localhead{\large\emph{#1}}}
\def\subheading#1{\localhead{\emph{#1}}}

\newenvironment{clist}%
{\bgroup\parskip 0pt\begin{list}{$\bullet$}{\partopsep 4pt\topsep 0pt\itemsep -2pt}}%
{\end{list}\egroup}%

\begin{document}

\strut\par\vskip-12pt

\fbox{
\begin{tabular}{lll}
{\emph{Instructor}}  & Leah Berman (\S 001) \\
& 
\href{mailto:lwberman@alaska.edu}{\texttt{lwberman@alaska\@@alaska.edu}} \cr
\strut & \cr
{\emph{Instructor Office}}&Chapman 102 \cr
\strut & \cr
{\emph{Office Hours}} & Announced on Canvas\cr
\strut & \cr
{\emph{Meeting Times}}& no set meeting time \cr
\strut & \cr
%{\emph{Classroom}}& Rasmuson Library& Rasmuson Library\\
% &Rm 660 & Rm 660 \\
%\strut &\\
{\emph{Prerequisites}} & \multicolumn{2}{l}{\parbox{4in}{ALEKS score $\geq$ 30, or a grade of B or better in MATH F055, MATH F062 or MATH F068, or permission of the instructor}}\\ 
\strut & \\
{\emph{Required text}} & \multicolumn{2}{l}{\textit{Math in Society} by David Lippmann}\cr
\strut & \cr
{\emph{Required Materials}} &  \multicolumn{2}{l}{\parbox{4in}{A non-programmable calculator\\
A printer or a tablet\\
A scanning device (a smartphone will work)}}\cr\\
{\emph{Course Website}} & \url{https://uaf-math.github.io/math113/} 
\strut and Canvas \cr
\end{tabular}
}
\strut
%\vspace{.1in}

\emph{\large{Overview \& Student Learning Outcomes}}

Mathematical reasoning shapes the decisions we make and the systems we navigate every day. This course explores how mathematical thinking helps us analyze problems ranging from the political to the personal. Topics will include the mathematics of voting systems; strategies for fair division; finding efficient routes and schedules; understanding financial mathematics; and introductory cryptography.

At the completion of the course, students will: 
	\begin{itemize}[noitemsep]
	\item know how to determine the winner of an election using a variety of voting methods and understand the strengths and weaknesses of each.
	\item understand a variety of strategies for equitably dividing things or tasks between parties.	
	\item use graphs to model commonplace, real-world applications and employ standard optimization algorithms.
	\item apply a variety of encryption methods and understand the strengths and weaknesses of each.
	\item understand some basic terminology and formulas of financial mathematics and explore their long-term consequences. 
	\end{itemize}

 This course is listed as a General Education Math Course. See GER Section at the end for details. 
 
\emph{\large{Information about this class}} In addition to this syllabus, information and materials for this class will be available on the Course Website \url{https://uaf-math.github.io/math113/} and in the Course Canvas class.

 \needspace{5\baselineskip}
\heading{Evaluation and Grades}
Grades are determined as follows.  (Each component of the grade is discussed below.)
 
\begin{multicols}{2}
\begin{tabular}{|c|c|}
\hline
%Participation & 8\%\\
%\hline
Homework & 8\% \\
\hline
Quizzes & 11\% \\
\hline
 Exams (3) & 3 $\times$ 24\% $=$ 72\% \\
\hline
Final Course Project & 9\% \\
\hline
{\bf total} & {\bf 100\%} \, \\
\hline
\end{tabular}
%
\hspace{1cm}
%\textcolor{red}{
\begin{tabular}{llll}
A  & 93--100\%& C  & 73--76\%  \\
A$-$ & 90--92\% & C$-$ & 70--72\% \\
B+ & 87--89\% & D+ & 67--69\%  \\
B  & 83--86\% & D  & 63--66\%  \\
B$-$& 80--82\% & D$-$ & 60--62\%  \\
C+ & 77-79\% & F  & $<$ 60\%
\end{tabular}
%}
\end{multicols}

%\heading{Participation in Classroom Activities}
%During many class periods, there will be a worksheet or short activity
%to be completed in class. These are cooperative assignments meant to
%be completed with your neighbors. The primary goal of the in-class work is to start the process of engaging with new ideas in the supportive environment of a classroom, making the outside-of-class learning --- via homework --- more productive. A student will only earn full participation points for actively participating in class activities.
%
%If you are not in class for a particular activity, you can still get participation points for absences by completing and submitting any worksheets you missed by the next class period. For more detailed instructions on how to make-up class participation points, consult with your instructor.
%
%Worksheets will be collected at the end of each class period and will be graded for completion and effort.
%
%
%\emph{Participation in Classroom Activities will make up 8\% of the overall course
%grade.} 

 \needspace{5\baselineskip}
\heading{Class Structure and Course Overview}
This is an asynchronous class, but it is \emph{not} a self-paced class. Each week, you are expected to learn certain class material, and you will need to submit homework assignments and quizzes (see below) on a regular basis as dictated by the course schedule, as well as take the exams at set times. 

Each week has its own Module in the Canvas Course. Each module will contain:
\begin{itemize}[noitemsep]
\item Lecture videos and associated notes templates for you to fill in as you watch the video
\item A link to the section in the textbook being covered
\item Worksheets that you should work through after you've watched each video, to help strengthen your understanding of the new material
\item A homework assignment, that you should work through after the videos and worksheets, and access to where you submit it
\item For weeks with quizzes, an access point into the quiz
\end{itemize}

Most weeks will have two or three lecture videos and associated worksheets. Solutions to the worksheets will be posted with the worksheet, and you are expected to check your answers against the solutions, and reach out for help if there are things you don't understand. The worksheets will not be collected. Homework will be submitted on an approximately weekly basis, and each week will either have a quiz or an exam.

There will be three exams, which will count for about 75\% of your grade; these exams \emph{must be proctored by a human} (more details below). You will want to start sorting out proctoring immediately so that it is set up by the time the exams come around; the first exam is during week 4 of the semester.

 \needspace{5\baselineskip}
\heading{Homework}
Throughout the semester, there will be weekly written homework assignments. Each assignment will be
due on Wednesday, as shown on the course schedule linked in Canvas/on the course website. You will turn them in as directed in Canvas. Homework problems are intended for practice. Complete worked solutions to the homework problems will be provided in advance so that you will be able to:
\begin{itemize}[nosep]
\item check your work
\item identify what you do not understand, 
\item \emph{fix any errors} you might have before you submit your homework
\item get help resolving any misunderstanding.
\end{itemize}
Effort and
completion will count towards a significant portion of your grade, but
accuracy does matter, as do organization and neatness. 

If you find it helpful to discuss homework problems with generative AI such as ChatGPT or Claude.ai, you are welcome to do so, but remember that the goal is that you develop your own understanding of the material. Just like simply copying the homework solutions is not likely to lead to understanding, copying the output of Claude or ChatGPT is also not likely to lead to understanding.

Homework problems must be written legibly and scanned and uploaded clearly (if your scan is unreadable you will not get credit for the assignment). See the \emph{Homework Guidelines} on the Homework page for detailed information about how your homework should be written.  


\emph{Homework is 8\% of the overall course
grade.}

 \needspace{5\baselineskip}
\heading{Quizzes}
At the end of most weeks, we will have a quiz over topics from the homework assignment that was due Wednesday. %Each quiz will be given in the last 30 minutes of class. You will have 15 minutes to complete the quiz independently, using only a nonprogrammable calculator. In the last 15 minutes of class, you will have the opportunity to correct your quiz with the help of your notes, your book, your classmates, and your instructor. Quizzes will be graded on correctness, but you will have the chance to get half of the points back for clearly identifying any mistakes and supplying detailed correct answers. 

Your quizzes are written quizzes with problems that will be similar to homework and exam problems.  They will be delivered to you through an online program; you will need to print out the quiz from the online program, work the quiz independently without access to course materials, scan the quiz, and upload the quiz back into the online program. (More instructions are provided in Canvas.) You will have 45 minutes to complete the downloading, quiz-taking, scanning, and uploading; the quizzes are designed to take you about 15 minutes to work the problems.

Your quizzes will not be proctored, so you are on your honor to not use outside materials, Generative AI programs, your classmates, or other course materials. I strongly encourage you to take the quizzes seriously as focused time to test your understanding of the material and to use as practice for the exams! You will have the opportunity to get additional points that you missed back from the quizzes (up to 50\% of missed points); more information will be provided on Canvas.

There will be an additional \emph{Finance Assessment} during the last week of class over the Finance material; this will count as two quizzes. More information will be distributed closer to the assessment time.

\emph{Quizzes are 9\% of the overall course
grade.}

 \needspace{5\baselineskip}
\heading{Exams}
There are 3 exams, scheduled during Week 4, Week 9, and Week 12. Specifically, the testing dates are as follows:
\begin{itemize}
\item \emph{Exam 1:} September 16 -- September 18
\item \emph{Exam 2:} October 21 -- October 23
\item \emph{Exam 3:} November 11 -- November 13
\end{itemize}

It is your responsibility to schedule your exam, either with UAF Testing Services or with another approved testing center, during these times. See \url{https://www.uaf.edu/ssc/testing/} for more information. They maintain a list of approved testing centers and can talk to you about the process of finding an approved proctor. If you have special circumstances, please send me an email or set up a meeting.

Each exam will be 1 hour. These exams are not cumulative; they will only test the material that has been covered since the previous exam. Non-programmable calculators will be allowed on these exams. Make-up exams will be given only for documented excused absences. If you know you have a conflict with an exam date due to a university-approved absence, please talk to your instructor as soon as possible.

\emph{Each  exam is 25\% of the overall course grade for a total of 75\% of the course grade.}

If a disaster hits and you miss an exam unexpectedly, email your instructor immediately!

 \needspace{5\baselineskip}
\subheading{Exam Corrections Policy}
For each of the three exams, students who score below 90\% will have the opportunity to correct their exams in order to learn from any mistakes. Moreover, a student may improve their exam grade in the process. A student can earn \emph{up to} 50\% of missed points for correcting their work, up to a maximum total score of 90\%. Only \emph{\underline{correct} new work} is eligible to receive additional points.
%\newpage
{\it For example,
\begin{itemize}
\item {\it A student who scores {\bf 50} out of 100 points may
correct all 50 points they missed and receive {\bf 25} points for
corrections, giving them a combined score of {\bf 75} points.}
\item {\it A student who scores {\bf 89} out of 100 points may
correct all 11 points they missed but will only receive {\bf 1} point
for corrections, capping the combined score at {\bf 90} points.}
\end{itemize}
}

Corrections are due one week after the graded exams are returned. They will be uploaded to the same online program where you take the quizzes. Detailed instructions are provided in Canvas.

 \needspace{5\baselineskip}
\heading{Final Course Project}
At the end of the semester, instead of taking a final exam, students will submit a Final
Course Project. Detailed instructions and standards will be posted in the \emph{Project Info} page linked in the left menu of the public webpage. A sample description of a project from Spring 2026 is linked there now. Each of the project options is designed to give the student a chance to apply some topic from the course in a setting of their choosing. The projects are due at the time of the final exam. 

\emph{The final course project is 9\% of the overall course grade}.

 \needspace{5\baselineskip}
\heading{Tutoring and Resources}
\vskip -30pt\strut
\begin{clist}
    \item The Math and Stat Lab is located in the 
      student success center on the 6th floor of the Rasmuson Library and
      offers drop-in tutoring, both in person and online!

	See 	\href{http://www.uaf.edu/dms/mathlab/}{\texttt{www.uaf.edu/dms/mathlab/}} for schedules and availability.
	\item Free
one-on-one (or small group) tutoring is also available. You must schedule an
appointment; see \href{http://www.uaf.edu/dms/mathlab/}{\texttt{www.uaf.edu/dms/mathlab/}}.
\item The Student Success Center has {\it Academic Coaches}, which are undergraduate students who can help you  improve your study strategies, identify resources and set goals, offer assistance with personalized study plans, time management,  navigating UAF technology, test and note-taking strategies, and much more. You can talk to Academic Coaches by dropping into their area in the Student Success Center on the 6th floor of the Rasmuson Library.
	\item Student Support Services (\href{https://uaf.edu/sss/}{\texttt{uaf.edu/sss/}}) offers free tutoring in many subjects to students who qualify for their program.
	\item ASUAF (\href{https://uaf.edu/asuaf/}{\texttt{uaf.edu/asuaf/}}) offers private tutoring for a small fee, based on student income.\\
\end{clist}

 \needspace{10\baselineskip}
\heading{Rules and Policies}

\subheading{Regular and Substantive Interaction}
To be compliant with federal law and Northwest Commission on Colleges and Universities (NWCCU) accreditation requirements, UAF faculty must ensure that all courses where faculty and students are not physically located in the same space for which students use federal financial aid have regular and substantive interaction (RSI) between students and instructors. Regular refers to interactions that are scheduled and predictable. Substantive refers to engaging students in teaching, learning, and assessment, consistent with course content. As your faculty member, I have included the following actions in the course to meet RSI requirements: 

%Providing direct instruction
\begin{itemize}[nosep]
\item Providing feedback on student coursework
\item Providing course-related information or answering questions
\item Facilitating group discussions about content or competencies
%Other instructional activities (please describe)
\end{itemize}

You will receive content-related feedback on all of your quizzes and exams, and in addition your homework assignments will be graded for completion. Of course, it is your responsiblity to review the feedback you get on your quizzes and exams and ask questions if there are things you don't understand!

Your canvas course shell has discussion boards set up for each major topic area, and you are encouraged to use these discussion boards to ask and answer questions. I will also monitor the discussion boards and also answer questions.

I will hold drop-in zoom office hours which you are encouraged to attend, and you also can set up individual appointments with me if those don't work with your schedule.

The regular and substantive feedback you receive will be assessed by your completion of a student survey; I will give \emph{extra credit} towards your homework grade if you complete the survey. The survey will open automatically around week 4 and will close around week 8. Further instructions on how to complete the survey will be sent to you directly from the survey committee.

 \needspace{5\baselineskip}
\subheading{AI usage policy} 
During proctored, paper exams, you will not have access to electronic tools of any type, and you may not use books or notes, except as announced.  These assessments represent around 75\% of your grade. 

You should not use Generative AI (or any other outside materials) to complete your quizzes.

Feel free to use a calculator or outside resources while completing your homework.  It is also reasonable to explore new AI tools like ChatGPT or Claude.ai. However, the point of this course is that you struggle with and learn and understand the material, not recite what a computer suggests, so it is important that you are using your own thinking when engaging with this course. In addition, since exams represent the  majority of your grade, as you do the homework, you must focus on your own thinking and level of understanding. Copying solutions without understanding will have no benefit to your own learning of the material, which is the goal of the homework.  Writing without understanding is also not a good long-term strategy for passing the course. 
  
Your final project \emph{must be completed using only your own words and ideas}. You may not use Generative AI to produce your final project. You are welcome and encouraged to talk to your classmates about your final project, but it must be completed individually.

\needspace{5\baselineskip}
\subheading{Incomplete Grade} 
An incomplete is a temporary grade used to indicate that the student has satisfactorily completed (C (2.0) or better) the majority of work in a course (usually all but the last 3 weeks) but for personal reasons beyond the student's control, such as sickness, has not been able to complete the course during the regular semester. See the catalog \url{https://catalog.uaf.edu/academics-regulations/grades/} for more details.

\subheading{Late Withdrawals} 
A withdrawal after the deadline from a DMS course will normally be granted only in cases where the student is performing satisfactorily (i.e., C or better) in a course, but has exceptional reasons, beyond their control, for being unable to complete the course. To apply for a late withdrawal, please talk to your instructor and your advisor; there is a separate university committee that evaluates such requests.

\subheading{Academic Dishonesty}
Academic dishonesty, including cheating and plagiarism, will not
be tolerated.  It is a violation of the Student Code of Conduct
and will be punished according to UAF procedures.

\subheading{Student protections and services statement}
Every qualified student is welcome in my classroom.  As needed, I am happy to work with you, Disability Services, Veterans' Services, Rural Student Services, etc.~to find reasonable accommodations.  Students at this University are protected against sexual harassment and discrimination (Title IX), and minors have additional protections.  As required, if I notice or am informed of certain types of misconduct, then I am required to report it to the appropriate authorities.  For more information on your rights as a student and the resources available to you to resolve problems, please go the following site: \url{https://catalog.uaf.edu/handbook/}.
 
\needspace{5\baselineskip}
\subheading{General education statement}
This course is listed as a General Education Math Course.  As such this course is expected to \textsl{contribute to meeting the following} four general learning outcomes:

\begin{enumerate}
\item Build knowledge of human institutions, sociocultural processes, and the physical and natural works through the study of mathematics.  Competence will be demonstrated for the foundational information in each subject area, its context and significance, and the methods used in advancing each.

\item Develop intellectual and practical skills across the curriculum, including inquiry and analysis, critical and creative thinking, problem solving, written and oral communication, information literacy, technological competence, and collaborative learning. Proficiency will be demonstrated across the curriculum through critical analysis of proffered information, well-reasoned solutions to problems or inferences drawn from evidence, effective written and oral communication, and satisfactory outcomes of group projects.

\item Acquire tools for effective civic engagement in local through global contexts, including ethical reasoning, intercultural competence, and knowledge of Alaska and Alaska issues.  Facility will be demonstrated through analyses of issues including dimensions of ethics, human and cultural diversity, conflicts and interdependencies, globalization, and sustainability.   

\item Integrate and apply learning, including synthesis and advanced accomplishment across general and specialized studies, adapting them to new settings, questions and responsibilities, and forming a foundation for lifelong learning. Preparation will be demonstrated though production of a a creative or scholarly product that requires broad knowledge, appropriate technical proficiency, information collection, synthesis, interpretation, presentation and reflection.
\end{enumerate}


\begin{center}
\heading{Official UAF Syllabus Addendum}
\end{center}


\noindent\emph{Student protections statement:} The university respects and upholds the principles of due process and a fair and equitable process as specified in the Board of Regents' Policy 09.02 Student Rights and Responsibilities. For more information regarding the rights and responsibilities of students, refer to the Office of Rights, Compliance and Accountability website. You are encouraged to read the Board of Regents' policy carefully to fully understand your responsibilities to our community.

We strive to create a safe and respectful environment for all members of our community. If you have questions about expectations of you as a student or believe your rights are being violated, we encourage you to reach out to the  Office of Rights, Compliance and Accountability for help. UAF reserves the right to suspend, expel or take other necessary and appropriate action in cases where a student is unable or unwilling to uphold community standards and campus safety.

For more information on your rights as a student and the resources available to you to resolve problems, please go to the following site:\\ \url{https://catalog.uaf.edu/academics-regulations/students-rights-responsibilities/}

\noindent\emph{Disability services statement:}  UAF is committed to creating a learning environment that meets the needs of its diverse student body.  If you have a disability, or think you might have a disability, please contact Disability Services (DS) to request a reasonable accommodation.  Accommodations are not retroactive, so students should request accommodations as early as possible, since they might take time to implement.  Students should notify DS at any time during the semester if adjustments to their communicated accommodation plan are needed, and must request their accommodations each semester.  More information is available online at \url{https://www.uaf.edu/disability}, and students can reach DS at uaf-disability-services@alaska.edu or 907-474-5655.


\noindent\emph{ASUAF Advocacy \& Engagement:} The Associated Students of the University of Alaska Fairbanks (ASUAF), UAF's student government, works to represent student voices and improve the student experience. Students are encouraged to get involved in shaping policies, sharing feedback, and helping address issues that impact campus life. ASUAF also offers support in navigating university processes and connecting with the right resources when challenges arise. To get involved or seek support, visit the ASUAF office or email asuaf.office@alaska.edu. 

 



\noindent\emph{Student Academic Support:}
\begin{itemize}
\setlength\itemsep{0em}
        \item Communication Center (907-474-7007, \mailto{uaf-commcenter@alaska.edu}, Student Success Center, 6th Floor Room 677 Rasmuson Library)
        \item Writing Center (907-474-5314, \mailto{uaf-writing-center@alaska.edu}, Student Success Center, 6th Floor Room 677 Rasmuson Library)
\item UAF Math Services (907-474-7332, \mailto{uaf-traccloud@alaska.edu})


\begin{itemize}
\item Drop-in tutoring, Student Success Center, 6th Floor Room 672 Rasmuson Library

\item 1:1 tutoring (by appointment only), 6th Floor Room 677 Rasmuson Library

\item Online tutoring (by appointment only)
available at the Student Success Center, schedule at \url{https://www.uaf.edu/dms/mathlab/}, 
\end{itemize}

\item Developmental Math Lab, Gruening 406
\item The Debbie Moses Learning Center at CTC (907-455-2860, 604 Barnette St, Room 102,\\ \url{https://www.ctc.uaf.edu/student-services/student-success-center/})
\end{itemize}

 For more information and resources, please see the Academic Advising Resource List (\url{https://www.uaf.edu/advising/students/index.php})


\noindent\emph{Student Resources:}
\begin{itemize}
\setlength\itemsep{0em}
\item Disability Services (907-474-5655, \mailto{uaf-disability-services@alaska.edu}, 110 Eielson Building)
\item Student Health \& Counseling [free counseling sessions available] (907-474-7043, \url{https://www.uaf.edu/chc/appointments.php}, Whitaker Building, Room 206, Health, Safety \& Security Bldg --- same building as Fire and Police)
\item Office of Rights, Compliance and Accountability (907-474-7300, \mailto{uaf-orca@alaska.edu}, 3rd Floor, Constitution Hall)
\item Associated Students of the University of Alaska Fairbanks (ASUAF) or ASUAF Student Government (907-474-7355, \mailto{asuaf.office@alaska.edu}{asuaf.office@alaska.edu}, Wood Center 119)
\end{itemize}

\noindent\emph{Nondiscrimination statement:}
The University of Alaska is an equal opportunity/equal access employer, educational institution and provider. The University of Alaska does not discriminate on the basis of race, religion, color, national origin, citizenship, age, sex, physical or mental disability, status as a protected veteran, marital status, changes in marital status, pregnancy, childbirth or related medical conditions, parenthood, sexual orientation, gender identity, political affiliation or belief, genetic information, or other legally protected status. The University's commitment to nondiscrimination, including against sex discrimination, applies to students, employees, and applicants for admission and employment. Contact information, applicable laws, and complaint procedures are included on UA's statement of nondiscrimination available at \url{www.alaska.edu/nondiscrimination}.


For more information, contact: 
\begin{tabular}{l}
UAF Office of Rights, Compliance and Accountability\\
1692 Tok Lane\\
3rd floor, Constitution Hall, Fairbanks, AK 99775\\
907-474-7300\\
\url{uaf-orca@alaska.edu}
\end{tabular}



\end{document}

%%% Local Variables:
%%% mode: LaTeX
%%% TeX-master: t
%%% End:
